% chapters/04/sub-cpmk.tex
\begin{subcpmkbox}
Bab ini menuntaskan \SUBCPMK{1.1} secara utuh. Ketiga indikatornya diperiksa
pada bab ini, dan buktinya berupa program siklus permainan yang Anda tulis
sendiri.

\vspace{0.5em}
\textbf{\SUBCPMK{1.1} Menjelaskan struktur \istilah{game loop} (input,
\istilah{update}, \istilah{render}) dan peran manajemen waktu dalam
permainan} --- ranah C2, \CPMK{1}, \CPL{03}.

\vspace{0.5em}
{\setlength{\tabcolsep}{4pt}\footnotesize%
\begin{tabularx}{\textwidth}{@{}C{1.9cm} L{4.0cm} L{2.5cm} Y@{}}
\toprule
\headrow \textbf{Indikator} & \textbf{Rumusan} & \textbf{Dibahas pada} & \textbf{Bukti} \\
\midrule
\ind{a} & Menjelaskan urutan tahapan
\istilah{input}--\istilah{update}--\istilah{render} dengan benar &
Bagian~\ref{sec:kerangka-siklus} &
Jawaban soal konsep dan penjelasan lisan atas program sendiri. \\
\addlinespace
\ind{b} & Mengidentifikasi peran \code{pygame.time.Clock} dan pembatasan
laju bingkai &
Bagian~\ref{sec:mengendalikan-laju} dan Bagian~\ref{sec:mengukur-laju} &
Program yang menampilkan laju bingkai pada judul jendela, beserta catatan
pengukuran pada lembar kerja praktikum. \\
\addlinespace
\ind{c} & Membedakan \istilah{fixed timestep} dan \istilah{variable timestep}
secara konseptual &
Bagian~\ref{sec:belum-ada} (diringkas dari Bagian~2.2) &
Jawaban soal uraian pada asesmen bab ini. \\
\bottomrule
\end{tabularx}}
\end{subcpmkbox}
